% ============================================
% Supplementary Note 6: Pool-size Dependency on Community-Level Selection
% ============================================
\subsection*{\rev{Supplementary Note 6: Pool-size dependency on community-level selection}}
\label{sec:suppl:pool_size_dependency}

\rev{We tested whether initial species-pool size could account for the observed community-level selection signatures in the experiments and simulations (Supplementary Fig.~34). In the synthetic-community experiments, we did not detect a significant effect of initial pool size on Dominance frequency across 6, 12, and 24 species (61.3\%, 58.9\%, and 59.4\%, respectively; $\chi^2 = 2.24$, $p = 0.69$), and this null result held within each nutrient condition (Dominance vs non-Dominance tests: Nutr$-$ $p = 0.39$, Base $p = 0.82$, Nutr$+$ $p = 0.27$). Realized parental richness increased with initial pool size (Kruskal--Wallis $p = 1.1 \times 10^{-6}$), while the ASV-based assembly survival ratio decreased with pool size (Kruskal--Wallis $p = 5.65 \times 10^{-13}$), but these assembly-richness effects did not translate into a corresponding pool-size dependence of Dominance frequency.}

\rev{In the gLV model, we performed a pool-size ablation at $\mu \in \{0.30, 0.60, 0.80\}$, varying the initial pool size from 4 to 48 species per parental community. Dominance frequency was primarily set by $\mu$ and changed little with pool size, whereas the same-parental-community versus cross-parental-community pairwise-selection gap increased with pool size. Reconstructed post-assembly interaction matrices showed lower within-community than between-community interaction strengths across pool sizes, supporting the interpretation that pool size affects assembly filtering and correlated species fates more than it affects Dominance-class frequency. Thus, pool-size variation does not explain the main transition in community-level selection; instead, it modulates the strength of correlated species fates within the interaction-strength regimes described in the main text.}
